\documentclass[11pt]{article}
\usepackage[margin=0.9in]{geometry}
\usepackage{amsmath,array,booktabs,enumitem,tikz}
\usepackage[version=4]{mhchem}
\usetikzlibrary{arrows.meta}
\setlength{\parindent}{0pt}\setlength{\parskip}{6pt}
\renewcommand{\arraystretch}{1.5}
\pagestyle{plain}
% orbital box: 0 empty, 1 up, 2 paired
\newcommand{\ob}[1]{\tikz[baseline=-0.7ex]{\draw (0,-0.28) rectangle (0.5,0.28);%
\ifcase#1\or\draw[-{Stealth[length=3pt]}](0.25,-0.2)--(0.25,0.2);%
\or\draw[-{Stealth[length=3pt]}](0.17,-0.2)--(0.17,0.2);\draw[-{Stealth[length=3pt]}](0.33,0.2)--(0.33,-0.2);\fi}}
\newcommand{\sub}[2]{\begin{tabular}[b]{@{}c@{}}#2\\[-6pt]\small #1\end{tabular}\hspace{10pt}}
\newcommand{\pb}[3]{\ob{#1}\hspace{-0.4pt}\ob{#2}\hspace{-0.4pt}\ob{#3}}
\newcommand{\ans}[1][1]{\par\foreach \i in {1,...,#1}{\vspace{0.8cm}\hrule}\vspace{4pt}}
\newcommand{\nm}{Name: \rule{7cm}{0.4pt}\par}

\begin{document}
{\Large\bfseries Study Guide: The Atom and the Periodic Table}\par\nm
The test on Monday has these same five parts.
\section*{Part 1: Rutherford's Gold Foil Experiment}
\textbf{1.} What did Rutherford fire at the gold foil, and what charge do they have?\ans

\textbf{2.} Fill in what each observation showed about the atom.

\begin{tabular}{|p{4.2cm}|c|p{7cm}|}\hline
What the alpha particles did & How many & What it showed\\\hline
Passed straight through & most & \\[12pt]\hline
Were deflected at an angle & a few & \\[12pt]\hline
Bounced straight back & very few & \\[12pt]\hline
\end{tabular}

\section*{Part 2: Inside the Atom}
\textbf{3.} Complete the table.

\begin{tabular}{|c|c|p{5cm}|}\hline
Particle & Charge & Where it is\\\hline
Proton & & \\\hline Neutron & & \\\hline Electron & & \\\hline
\end{tabular}

\textbf{4.} Atomic number ($Z$) = the number of \rule{3cm}{0.4pt}. \quad Mass number ($A$) = \rule{2.5cm}{0.4pt} + \rule{2.5cm}{0.4pt}.

\textbf{5.} Complete the table.

\begin{tabular}{|c|c|c|c|c|c|}\hline
Element & $Z$ & $A$ & Protons & Neutrons & Electrons\\\hline
O & 8 & 16 & & & \\\hline
Al & & 27 & & & \\\hline
 & & & 20 & 20 & \\\hline
\end{tabular}

\textbf{6.} What is an isotope? Give one example pair.\ans[2]

\textbf{7.} Boron is 20\% B-10 and 80\% B-11. Calculate its relative atomic mass.\ans[2]

\section*{Part 3: Electrons}
\textbf{8.} Write the filling order up to $4s$. \rule{8cm}{0.4pt}

\textbf{9.} Write the electron configuration for phosphorus (15 electrons) and for argon (18 electrons).\ans[2]

\textbf{10.} Fill in the orbital boxes for fluorine (9 electrons).\par\vspace{6pt}
\sub{1s}{\ob0}\sub{2s}{\ob0}\sub{2p}{\pb000}

\textbf{11.} For phosphorus: energy levels \rule{1.5cm}{0.4pt} \quad core electrons \rule{1.5cm}{0.4pt} \quad valence electrons \rule{1.5cm}{0.4pt}

\section*{Part 4: Ions}
\textbf{12.} A cation has \rule{2cm}{0.4pt} electrons than protons and a \rule{2cm}{0.4pt} charge. An anion has \rule{2cm}{0.4pt} electrons than protons and a \rule{2cm}{0.4pt} charge.

\textbf{13.} Complete the table.

\begin{tabular}{|c|c|c|c|}\hline
Ion & Protons & Electrons & Cation or anion?\\\hline
\ce{K+} & & & \\\hline \ce{S^2-} & & & \\\hline \ce{Ca^2+} & & & \\\hline \ce{F-} & & & \\\hline
\end{tabular}

\textbf{14.} A lithium atom becomes \ce{Li+}. Explain what happens to its electrons and why the charge is $1+$.\ans[2]

\textbf{15.} Which ion has the same electron configuration as neon: \ce{K+}, \ce{Na+} or \ce{Cl-}?\ans

\section*{Part 5: The Periodic Table and Trends}
\textbf{16.} Complete the table.

\begin{tabular}{|c|p{3.5cm}|c|c|p{4cm}|}\hline
Group & Name & Valence $e^-$ & Ion & Reactivity\\\hline
1 & & & & \\\hline 2 & & & & \\\hline 17 & & & & \\\hline 18 & & & & \\\hline
\end{tabular}

\textbf{17.} What is first ionization energy?\ans

\textbf{18.} Which has the higher ionization energy, magnesium or sulfur? Why?\ans[2]

\textbf{19.} Which atom is larger, sodium or rubidium? Why?\ans[2]

\textbf{20.} Why are noble gases so unreactive?\ans
\end{document}
